\section{相关工作}

\paragraph{面向 Agent 的长期记忆。} 近期工作把 Agent 记忆当作一种架构，而不是一个向量库。
ZenBrain 提出一个七层、受神经科学启发的记忆架构，包含睡眠期巩固（sleep-phase
consolidation）与对每个机制的显式消融~\cite{zenbrain2026}。FSFM 主张选择性遗忘应被当作一等问题对待，
并给出被动衰减、主动删除、安全触发遗忘与自适应强化四类分类~\cite{fsfm2026}。
一个受人类启发的架构结合睡眠期巩固、基于干扰的遗忘与多线索检索，
在存储缩减 58\% 的情况下报告 97.2\% 的保留精度（retention precision）~\cite{himemory2026}。
PersMem 把人格内化进一条双路径记忆流水线，使记忆处理变得依赖人格~\cite{persmem2026}。
这些系统优化的是「记住什么」与「检索得多好」。我们的关切不同且正交：
什么被\emph{允许}进入 prompt、按什么优先级进入，以及当预算装不下全部内容时会发生什么。

\paragraph{激活与衰减。} 在长时间跨度上对记忆排序，通常被表述为衰减或激活问题。
ACT-R 基础级激活（base-level activation）把频率与近因折叠成一个对数量，
其动机来自练习与需要概率的统计规律~\cite{anderson1991,anderson2007}。
我们把它当作两条通道之一，而不是衰减的替代品；因为在我们的设定里，
彻底替换会破坏一个冻结的排序基线，这一点报告在第 4 节。
这是一个普遍观点的具体例子：借来的机制很少能在不经受「与借用之前已存在的工具对照的回归检查」的情况下移植成功。

\paragraph{情感与人格。} Sentipolis 把连续的 Pleasure--Arousal--Dominance 表示与双速情感动力学、
以及情感—记忆耦合结合起来，用于社会模拟~\cite{sentipolis2026}。它自己的报告很有启发：
收益依赖模型，而且情感感知可能轻微降低对社会规范的遵守。PAD 本身可追溯到环境心理学~\cite{mehrabian1974}。
我们采纳 PAD 作为表示，但与面向模拟的系统不同，把情感限制在风格域：
一条不变量禁止情感进入事实判断，而合成器在结构上强制这种分离。

\paragraph{安全工程。} STPA 分析安全关键系统中的不安全控制动作（unsafe control
actions）~\cite{leveson2011,leveson2018}。我们把「不安全控制动作」这一概念从它原本离线的
分析场景，改造为一道运行时闸门：八个控制动作在组装之后、模型调用之前被筛查，
并且只有一个 fail-closed 落点。这一移植是刻意的，但只是部分的——STPA 是一种危险分析方法，
而我们的闸门是每一轮都要评估的谓词——并且我们不声称这道闸门可以替代对宿主做一次完整的 STPA。

\paragraph{控制论与必要多样性。} Ashby 的必要多样性定律指出，
调节器所掌握的多样性至少要与它面对的扰动一样多~\cite{ashby1956}；
Beer 的可生存系统模型（Viable System Model）是关于分层调节的一种有影响力的读法~\cite{beer1972}。
我们借用一个工程想法：当控制器无法表示某个冲突时，它应当收缩自己的动作空间，而不是照旧行动；
我们把它实现为跨合成域的降维。我们只把 VSM 当作对一套既有优先级分层的解释性透镜：
它是被借用的，不是被实现的，在我们的实现里没有代码落点、参数或依赖，
我们也不宣称 System~1--5 的机制。积分控制器的抗饱和是标准控制实践~\cite{astrom2006}，我们直接采纳。

\paragraph{评估与分阶段变更。} 只优化单一代理指标会招致 Goodhart 定律~\cite{goodhart1975}，
因此我们的学习回路把每个主指标与一个对抗项配对，把用户纠正当作最终裁判，
并冻结评估工具本身。变更管理遵循常规的分阶段放量模式（小比例、护栏、自动回滚），
并附加一条约束：不足以验证某条护栏的数据会阻断发布，而不是跳过检查。

\paragraph{范围。} 上文讨论的机制构成本文范围的全部。该范围的边界——
包括项目中存在但在这里刻意不做声称的组件——在第 7 节说明。
